  \documentclass[11pt,a4paper]{article}
\usepackage[utf8]{inputenc}
\usepackage[T1]{fontenc}
\usepackage{geometry}
\usepackage{parskip}
\usepackage{hyperref}
\geometry{margin=2.5cm}
\hypersetup{colorlinks=true,urlcolor=blue!60!black}

\pagestyle{empty}

\begin{document}

\begin{flushright}
\textbf{Fouad Bensmail}\\
Independent researcher\\
Constantine, Algeria\\
ORCID: 0009-0006-9541-4207\\
Email: \href{mailto:bensmfouad@gmail.com}{bensmfouad@gmail.com}\\[0.5em]
14 September 2026
\end{flushright}

\vspace{1em}

\noindent To the Editors,\\
\textit{Experimental Mathematics}\\
Taylor \& Francis

\vspace{1.5em}

\noindent \textbf{Subject:} Submission of ``Weighing the obstruction: the smoothness of $p-1$, its parity signature, and a measured Tamagawa reconstruction of the twin-prime constant''

\vspace{1em}

\noindent Dear Editors,

I am pleased to submit my manuscript entitled \emph{Weighing the obstruction: the smoothness of $p-1$, its parity signature, and a measured Tamagawa reconstruction of the twin-prime constant} for consideration as a research article in \emph{Experimental Mathematics}.

This work is experimental in the exact spirit of the journal's manifesto: every claim is a high-precision measurement, reproducible by a single dependency-free Python command, and each measurement carries its own guard (a predicted ratio, a sigma budget, and a contraction test across two scales).

The paper reports four main results:

\begin{enumerate}
\item The naive Dickman model for $P^{+}(p-1)$ is rejected at $10^{6}$. A corrected null model (parity plus Dirichlet densities $1/(q-1)$) is then forged, and its residual evaporates as $1/\sqrt{N}$ across two scales, excluding a $1/\ln X$ bias.

\item Conditioning on $p \bmod 4$, the local obstruction at $q=2$ proves computable and is captured class by class (8/8 ratios within $2\sigma$), while the unconditioned model misses each class in the predicted direction (structural bias, not fluctuation).

\item The singular series of pair motifs is proven to be a relative invariant: exact under translations and reflections, covariant under dilation by a multiplicative character verified to twelve digits. This covariance descends into prime-pair counts, reproducing the published twin counts ($58\,980$ at $10^{7}$, $440\,312$ at $10^{8}$) exactly.

\item Chinese Remainder independence is measured rather than assumed, and local Tamagawa volumes $\beta_{p}=1-1/(p-1)^{2}$ reconstruct the twin-prime constant $2C_{2}=1.3203236$ to $1.2\times 10^{-7}$.
\end{enumerate}

\textbf{Honesty of scope.} The paper proves no open conjecture and claims none. Four of the five architectural layers are unconditional theorems (relative invariance; CRT with explicit error; Euler tail; cross-term extinction at the integer layer via smooth numbers in fixed-modulus progressions, after Tenenbaum, Chap.~III.5). The fifth layer (the prime layer) is stated as a measured conjecture, whose proof is guarded by the Hardy--Littlewood family itself. Furthermore, two initial predictions were corrected by the measurement itself, and one self-detected coding error is published as a refused run in the companion log.

\textbf{Reproducibility and Data Availability.} All companion scripts (pure Python, no dependencies) and the full laboratory notebook are archived on Zenodo. The preprint of this manuscript is available at DOI \href{https://doi.org/10.5281/zenodo.22758846}{\texttt{10.5281/zenodo.22758846}}, and the laboratory notebook is archived under concept DOI \href{https://doi.org/10.5281/zenodo.22314057}{\texttt{10.5281/zenodo.22314057}} (cahier v9: \href{https://doi.org/10.5281/zenodo.22757396}{\texttt{10.5281/zenodo.22757396}}; v10 contains the formal proofs T1--T4). Every table in the manuscript is regenerable by the single commands listed in Section~8.

This manuscript is original, has not been published previously, and is not currently under consideration for publication elsewhere. As the sole author, I confirm that I have no conflicts of interest to disclose and that this research received no specific grant from any funding agency.

Thank you for your time and consideration. I look forward to your feedback.

\vspace{2em}

\noindent Sincerely,

\vspace{2em}

\noindent \textbf{Fouad Bensmail}\\
Independent researcher, Constantine, Algeria\\
ORCID: \href{https://orcid.org/0009-0006-9541-4207}{0009-0006-9541-4207}

\end{document}